\part*{Preámbulo}

Tomé estos apuntes cuando cursé Diseño y Análisis de Algoritmos en la Universidad de los Andes, entre agosto y noviembre de 2022. 

Posteriormente los fui corrigiendo y extendiendo de a ratos, cuando trabajaba con temas de algoritmia. Las extensiones más notables corresponden a dos momentos: cuando cursé Teoría de Grafos entre agosto y noviembre de 2023, y cuando fui seleccionado para dictar el curso de Diseño de Algoritmos en mayo de 2026.

Los apuntes están presentados en el mismo formato en el que originalmente fueron escritos: la clase de \LaTeX{} que establece su estilo la escribí en noviembre del 2020. Más importante aún, los apuntes preservan la estructura y el tono de cuando los escribí como estudiante, con la intención de que sean útiles precisamente a esa audiencia.

% TODO: Ojo!
Espero que le sean útiles a todo aquel quien, como yo, deteste la notación que se usa en el libro de Cormen.

% TODO: Aclaraciones: quién es la audiencia del CLRS, quién es mi audiencia?
% Dónde se evidencia eso?
% En merge sort, prefiero slicing que pasar idxs por todo lado
% En max subarreglo con divide conquer, CLRS pone if-else menos legibles pq quieren algo más parecido a una prueba matemática; en SW de producción, uno querría un early return